

\section{Fixed points on affine Grassmannians}\label{section: fixed points on affine grassmannians}

Let $G$ be a reductive group and $\Gr_G = \lo G / \lop G$ its affine Grassmannian over $k$. We assume that $k$ is algebraically closed field. For ease of discussion, we further assume $k$ is of characteristic zero; this is not essential and can be largely avoided.

Consider the fixed-point family $\Fix_{\lo G \rtimes \Gm^{\rot}}(\Gr_G)$. Since $\Gr$ is an ind-scheme, so is $\Fix_{\lo G \rtimes \Gm^{\rot}}(\Gr)$.
Its natural derived structure is trivial by the criterion from \S \ref{subsubsection: reduced and derived variants} applied to the ind-presentation by affine Schubert varieties.


Our main interest lies in the classical reduced restriction of this family to the extended torus
\begin{equation*}
    \Fix^{\red}_{T \times \Gm^{\rot}}(\Gr_G).
\end{equation*}
In the rest of this section, we explain how to effectively describe the fibers of this family as the point $(x, \zeta) \in T \times \Gm^{\rot}$ varies, for both $\Gr_G$ as well as for affine Schubert varieties $\Gr_{\leq \mu}$ inside it.
%
We separate the section into a uniform discussion in \S \ref{section: general structure of the fixed points}, followed by specific descriptions in the case of $\zeta$ generic \S \ref{section: zeta generic} and $\zeta$ root of unity \ref{section: zeta root of unity}.





\subsection{General structure of the fixed points}\label{section: general structure of the fixed points}

\subsubsection{Centralizer action on fixed-points.}\label{section: the subgroup H^+}
With the subgroups $J \leq \lo G$ and $J^+ \leq \lop G$ from \S \ref{section: centralizers}, the following statements are obvious.
\begin{lemma}\label{lemma: H acts on fixed points general case}
The subgroup $J$ acts on $\Fix_{(x, \zeta)}(\Gr_G)$.
\end{lemma}
\begin{proof}
Since $J \leq \Cent_{\lo G \rtimes \Gm^{\rot}}(x, \zeta)$ is pointwise fixed under the adjoint action of $(x, \zeta)$ on $\lo G \rtimes \Gm^{\rot}$, it acts on the corresponding fiber of the fixed-point scheme by \S \ref{section: residual action on Fix}.
\end{proof}
%
\begin{lemma}\label{lemma: H^+ acts on fixed points general case}
The subgroup $J^{+}$ acts on $\Fix_{(x, \zeta)}(\Gr_{\leq \mu})$. 
\end{lemma}
\begin{proof} 
Since $\lop G\rtimes \Gm^{\rot}$ acts on $\Gr_{\leq \mu}$ and $J^+ = J \cap \lop G$, this follows by restricting the action from Lemma \ref{lemma: H acts on fixed points general case}.
\end{proof}

\begin{remark}
In the above statements, note that the whole centralizer $\Cent_{\lo G \rtimes \Gm^{\rot}}(x, \zeta)$ and its positive part $\Cent_{\lop G \rtimes \Gm^{\rot}}(x, \zeta)$ act on the respective fixed points. This amounts to an extra action by $\Gm^{\rot}$ as well as by the group of connected components. This is a natural piece of structure, but we do not need it for our purposes.
\end{remark}







\subsubsection{Fixed points on affine Grassmannian.}
We first describe the fixed points on $\Gr_G$ in terms of the connected centralizer $J$. 

To this end, recall the reflection subgroup $W_J \leq W_{\aff}$ from Notation \ref{notation: WJ}. We then have the labeling set $\Theta = X_{\bullet}/W_J$ from Notation \ref{notation: Theta}. 
%
For any $\vartheta \in \Theta$, we may choose a lift $\lambda \in X_{\bullet}$ and put $J^+_{\vartheta} := J^+_{\lambda}$ in the notation \eqref{equation: positive stabilizer}. This is independent of the choice of $\lambda$ up to isomorphism by Remark \ref{remark: independece of J vartheta on the lift} below.


\begin{lemma}\label{lemma: fixed points on affine Grassmannian}
 Let $(x, \zeta) \in T \times \Gm^{\rot}$. Then
 \begin{equation}\label{equation: fixed points on Gr}
     \Fix^{\red}_{(x, \zeta)}(\Gr_G) = \coprod_{\vartheta \in \Theta} J/J^{+}_{\vartheta}. 
 \end{equation}
\end{lemma}
\begin{proof}
Any point of $x \in \Gr_G$ can be written as a product
\begin{equation}\label{equation: root decomposition in affine grassmannian}
t^{\lambda} \cdot u_1 \cdots u_d
\end{equation}
with $\lambda \in X_{\bullet}$ and $u_i \in U_{\beta_i + m_i\hbar}$ nontrivial points in finitely many pairwise distinct affine root groups; see \cite[proof of Lemma 4.4]{RW22} for this well-known fact (\cite{RW22} discusses a more refined statement for a given Iwahori orbit -- our variant follows since the Iwahori orbits of $t^{\lambda}, \lambda \in X_{\bullet}$ are affine spaces paving the affine Grassmannian). 

The action of $(x, \zeta)$ fixes each $t^{\lambda}$. We now restrict attention to the fixed points in \eqref{equation: root decomposition in affine grassmannian} for one chosen $t^{\lambda}$. The expression \eqref{equation: root decomposition in affine grassmannian} is unique. On each root subgroup $U_{\beta + m\hbar}$, the element $(x, \zeta)$ acts via $\beta(x) \cdot \zeta^m$. Comparing to Lemma \ref{lemma: root spaces of J}, we deduce that the fixed points of this form are given by the orbit of $J$ through $t^{\lambda}$. Since $J^+_{\lambda}$ is the stabilizer of $t^{\lambda}$, this orbit is given by $J/J^+_{\lambda}$.

Ranging over all $\lambda \in X_{\bullet}$, we see that the fixed points decompose into a disjoint union of the orbits $J/J^+_{\lambda}$ of $t^{\lambda}$, $\lambda \in X_{\bullet}$. Moreover, two points $t^{\lambda_1}, t^{\lambda_2}$ lie in the same orbit of $J$ if and only if $\lambda_1, \lambda_2$ can be connected by affine reflections from $\Phi^{\bullet}_J$, in other words if and only if $\lambda_1, \lambda_2$ give the same point $\vartheta \in \Theta = X_{\bullet} / W_{J}$.
The lemma follows.
\end{proof}

\begin{remark}\label{remark: independece of J vartheta on the lift}
The final paragraph of the above proof shows that $t^{\lambda_1}$, $t^{\lambda_2}$ lie in the same orbit of $J$ if and only if $\lambda_1, \lambda_2$ give the same point $\vartheta \in \Theta = X_{\bullet} / W_{J}$. In particular, the stabilizer $J^{+}_{\lambda}$ depends, up to isomorphism given by conjugation inside $J$, only on the coset $\vartheta \in \Theta$ of $\lambda$. It is thus fair to denote it by $J^+_{\vartheta}$.
\end{remark}


\begin{notation}
We denote the connected components of \eqref{equation: fixed points on Gr} by
\begin{equation*}
    \Fix^{\red, \vartheta}_{(x, \zeta)}(\Gr_G) = J/J^{+}_{\vartheta}, \qquad \vartheta \in \Theta.
\end{equation*}
\end{notation}




\subsubsection{Fixed points on $\Gr_G$ and $\Gr_L$.}\label{section: fixed points on GrG and GrL}
We also have the following alternative approach for describing the fixed points $\Fix^{\red}_{(x, \zeta)}(\Gr_G)$.
%
Consider the defining inclusion of the resonant subgroup $L \leq G$. It induces a map on loop groups and on affine Grassmannians $\Gr_L \subseteq \Gr_G$. Using this embedding, we obtain the following identification of fixed points.

\begin{lemma}\label{lemma: replacing x by the resonant part}
We have identifications
\begin{equation*}
    \Fix^{\red}_{(x, \zeta)}(\Gr_G) = \Fix^{\red}_{(\omega(\zeta), \zeta)}(\Gr_L).
\end{equation*}
\end{lemma}
\begin{proof}
Consider the natural inclusion $\Gr_L \subseteq \Gr_G$ induced from $L \leq G$. Both $L$ and $G$ have the common maximal torus $T$. Both affine Grassmannians thus have the same coordinate fixed points $t^{\lambda}$, $\lambda \in X_{\bullet}$. Fixing one such $t^{\lambda}$ and using the decomposition \eqref{equation: root decomposition in affine grassmannian}, we are left to compare which root spaces are fixed under the adjoint action of $(x, \zeta)$. A root space $U_{\beta + m\hbar}$ of $\lo G$ for $\beta + m\hbar \in \Phi^{\bullet}_{G, \aff}$ is fixed if and only if $\beta(x)\cdot \zeta^m=1$. For this to be true, it is necessary that $\beta \in \Phi^{\bullet}_L$. For such $\beta \in \Phi^{\bullet}_L$, we have $\beta(x)=\beta(\omega(\zeta))$ by construction of $\omega$. Hence the above condition is satisfied if and only if $\beta \in \Phi^{\bullet}_L$ and $\beta(\omega(\zeta))\cdot \zeta^m=1$. This precisely gives the root spaces of $\lo L$ fixed under the adjoint action of $(\omega(\zeta), \zeta)$ as desired.
\end{proof}

Moreover, we can get rid of the element $\omega(\zeta)$ at the cost of a shift. 
To this end, take the adjoint form $L^{\ad}$ of $L$. Since any connected component of $\Gr_L$ is equivariantly isomorphic to the corresponding connected component of $\Gr_{L^{\ad}}$ via the map $\pi_1(L) \to \pi_1(L^{\ad})$ of their labeling sets, we do not lose anything by working with $L^{\ad}$. Here $\omega$ is a genuine cocharacter, so multiplication by $t^{\omega} \in \lo L^{\ad}$ induces a shift map on the corresponding affine Grassmannian $\Gr_{L^{\ad}}$.

\begin{lemma}\label{lemma: fixed points and the omega shift general case}
Translation by $t^{\omega}$ induces an eqivariant isomorphsim 
\begin{equation}
t^{\omega} \cdot \Fix_{(\omega(\zeta), \zeta)}(\Gr_{L^{\ad}}) = \Fix_{(1, \zeta)}(t^{\omega} \cdot \Gr_{L^{\ad}}).   
\end{equation}
\end{lemma}
\begin{proof}
Let $p \in \Gr_{L^{\ad}}$. We need to see that $(\omega(\zeta), \zeta)\cdot p = p$ holds if and only if $(1, \zeta) \cdot \omega(t) \cdot p = \omega(t) \cdot p$. By the general structure of fixed-point schemes \S \ref{section: residual action on Fix}, it is enough to check that inside $\lo L^{\ad} \rtimes \Gm^{\rot}$ we have
\begin{equation}
 (\omega(\zeta), \zeta) = \omega^{-1}(t) \cdot (1, \zeta) \cdot \omega(t).
\end{equation}
This is the case by direct computation within $\lo L^{\ad} \rtimes \Gm^{\rot}$:
\begin{align*}
(\omega(\zeta), \zeta)
&= (\omega(t^{-1}) \cdot 1 \cdot \omega(\zeta t), 1 \cdot \zeta \cdot 1)\\
&= \cdot (\omega(t^{-1}), 1) \cdot (1, \zeta) \cdot (\omega(t), 1)\\
&= \omega^{-1}(t) \cdot (1, \zeta) \cdot \omega(t).
\end{align*}
where the second equality comes from the semidirect product $\lo L^{\ad} \rtimes \Gm^{\rot}$ since conjugation by the loop rotation factor scales $t$ with weight one.
\end{proof}



Putting Lemmata \ref{lemma: replacing x by the resonant part}, \ref{lemma: fixed points and the omega shift general case} together, we have identified $\Fix^{\red}_{(x, \zeta)}(\Gr_G)$ as a disjoint union of fixed points of the loop rotation by $\zeta$ on a disjoint union of connected components $\Gr_{L^{\ad}}^{\sigma}$ of the affine Grassmannian of $L^{\ad}$

\begin{equation}\label{equation: fixed points on Gr_G and loop rotation on Gr_L}
    \Fix^{\red}_{(x, \zeta)}(\Gr_G) 
    \cong \coprod_{\pi_1(L)} \Fix^{\red}_{(1, \zeta)}(t^{\omega} \cdot \Gr^{{\sigma}}_{L^{\ad}})
    = \coprod_{\pi_1(L)} \Fix^{\red}_{(1, \zeta)}(\Gr^{{\underline{\omega} + \sigma}}_{L^{\ad}})
\end{equation}
The disjoint union runs over the set $\pi_1(L)$ of connected components of $\Gr_L$, whose elements determine $\sigma \in \pi_1(L^{\ad})$ thought the natural map $\pi_1(L) \to \pi_1(L^{\ad})$.


The fixed points of the loop rotation by $\zeta$ on affine Grassmannians -- such as $\Gr_{L^{\ad}}$ -- are well known. The identification \eqref{equation: fixed points on Gr_G and loop rotation on Gr_L} thus reduces the computation of $\Fix^{\red}_{(x, \zeta)}(\Gr_G)$ to this known special case.






\subsubsection{Fixed-points on affine Schubert varieties.}
We can further refine the computation from Lemma \ref{lemma: fixed points on affine Grassmannian} to describe fixed points on each affine Schubert variety $\Gr_{\leq \mu}$ in $\Gr_G$.
From \S \ref{section: the subgroup H^+} we get the following diagram
\begin{equation}\label{equation: fixed points on affine schubert varieties 1}
J^+ \acts
\begin{tikzcd}
    \Gr_{\leq \mu}  \arrow[r, hookrightarrow] & \Gr_G \arrow[r, equal] & \lo G / \lop G \\
    \Fix_{(x, \zeta)}(\Gr_{\leq \mu}) \arrow[u, hookrightarrow] \arrow[r, hookrightarrow] & \Fix_{(x, \zeta)}(\Gr_G) \arrow[u, hookrightarrow]  \arrow[r, equal] & \coprod_{\vartheta \in \Theta} J/J^{+}_{\vartheta}. \arrow[u, hookrightarrow]
\end{tikzcd}
\end{equation}
where all maps are equivariant closed immersions.

\begin{lemma}\label{lemma: upshot of the general case}
The scheme $\Fix^{\red}_{(x, \zeta)}(\Gr_{\leq \mu})$ identifies with a disjoint union of closed $J^+$-equivariant subschemes
\begin{equation*}\label{equation: B-equivariant closed immersion into classical flag varieties general case}
J^+ \acts \Fix^{\red}_{(x, \zeta)}(\Gr_{\leq \mu}) \hookrightarrow \coprod_{\vartheta \in \Theta} J/J^+_{\vartheta}.
\end{equation*}
\end{lemma}
\begin{proof}
Take the reduced bottom row of diagram \eqref{equation: fixed points on affine schubert varieties 1}. 
\end{proof}


Furthermore, it is useful to understand the embedding from Lemma \ref{lemma: upshot of the general case} in terms of the coordinate fixed points $t^{\lambda}$. Given any $\mu \in X^{+}_{\bullet}$, recall \eqref{notation: weights in mu} that 
$ \wt(\mu) = W \cdot \{ \lambda \in X^{+}_{\bullet} \mid \lambda \leq \mu \} \subseteq X_{\bullet}$
stands for the union of the Weyl orbits of the dominant weights $\lambda$ below $\mu$.
Under the geometric Satake equivalence, $\wt(\mu)$ labels the fixed points of $T \times \Gm^{\rot}$ on $\Gr_{\leq \mu}$.
Consequently, the fixed point scheme $\Fix_{(x, \zeta)}(\Gr_{\leq \mu})$ contains the discrete set $t^{\lambda}$, $\lambda \in \wt(\mu)$. 


\begin{lemma}\label{lemma: connected components of fixed points n affine schubert variety}
The $T\times \Gm^{\rot}$-fixed-points on $\Fix_{(x, \zeta)}(\Gr_{\leq \mu})$ are given by $\{ t^{\lambda} \mid \lambda \in \wt(\mu) \}$. Two such fixed-points $t^{\lambda_1}$, $t^{\lambda_2}$ lie in the same connected component of $\Fix_{(x, \zeta)}(\Gr_{\leq \mu})$ if and only if $\lambda_1$, $\lambda_2$ lie in the same orbit of $W_J \acts X_{\bullet}$.   
\end{lemma}
\begin{proof}
We may work with reduced fixed point schemes.
The description of coordinate fixed points is clear from the discussion above. We are left to describe the connected components.


For one direction, assume that $\lambda_1, \lambda_2$ lie in the same orbit of $W_J$. Then they can be connected by a sequence of affine reflections $s_{\beta + m\hbar}$ for $\beta+m\hbar \in \Phi^{\bullet}_J \subseteq \Phi^{\bullet}_{G, \aff}$. To see this, first assume that $\lambda_2 = s_{\beta + m\hbar} \cdot \lambda_1$ for a single reflection in $W_J$. We claim that the affine root $\beta + m\hbar$ induces a $\P^1$ which connects the points $t^{\lambda_1}$, $t^{\lambda_2}$ inside $\Fix^{\red}_{(x, \zeta)}(\Gr_{\leq \mu})$.
%
Indeed, whenever two coweights $\lambda_1$, $\lambda_2$ are connected by an affine reflection $s_{\beta+m\hbar} \in W_{G, \aff}$, then the associated root space $U_{\beta + m\hbar}$ extends to a $\P^1$ connecting the corresponding points $t^{\lambda_1}$, $t^{\lambda_2}$ inside $\Gr_G$. Moreover, if $\lambda_1, \lambda_2 \in \wt(\mu)$, this $\P^1$ lies inside $\Gr_{\leq \mu}$. Since we also have ${\beta + m\hbar} \in \Phi^{\bullet}_J$, the root space $U_{\beta + m\hbar}$ lies inside $J$, hence the induced $\P^1$ actually lies inside $\Fix^{\red}_{(x, \zeta)}(\Gr_G)$. Altogether, it lies inside 
\begin{equation*}
    \Fix^{\red}_{(x, \zeta)}(\Gr_{\leq \mu}) = \Fix^{\red}_{(x, \zeta)}(\Gr_G) \cap \Gr_{\leq \mu},
\end{equation*}
proving the claim. 
%
In general, choosing a sequence of affine reflection of $W_J$ connecting $\lambda_1$, $\lambda_2$, we have thus constructed a chain of $\P^1$'s connecting $t^{\lambda_1}$ and $t^{\lambda_2}$ inside $\Fix^{\red}_{(x, \zeta)}(\Gr_{\leq \mu})$. Hence they lie in the same connected component.


For the converse, note that when $\lambda_1$ and $\lambda_2$ lie in different orbits of $W_J$, the fixed-points  $t^{\lambda_1}$, $t^{\lambda_2}$ end up in different terms of the disjoint union $\coprod_{\vartheta \in \Theta} J/J^+_{\vartheta}$, hence lie in different connected components of its subscheme $\Fix^{\red}_{(x, \zeta)}(\Gr_{\leq \mu})$.
\end{proof}

To conclude our computation of the fixed points, it remains to explicitly describe the players in Lemmata \ref{lemma: fixed points on affine Grassmannian}, \ref{lemma: upshot of the general case}. We do so in the next two subsections, distinguishing whether $\zeta$ is generic or a root of unity.








\subsection{\texorpdfstring{$\zeta$}{Zeta} generic}\label{section: zeta generic}
\subsubsection{Setup.}
Let $(x, \zeta) \in T \times \Gm^{\rot}$ be a closed point with $\zeta$ generic. The purpose of this subsection is to describe the fibers of the fixed-point scheme $\Fix^{\red}_{(x, \zeta)}(\Gr_G)$ as infinite disjoint unions of partial flag varieties for the resonant subgroup $L = L_{[x, \zeta]}$, as well as $\Fix^{\red}_{(x, \zeta)}(\Gr_{\leq \mu})$ as a union of finite Schubert varieties inside them.

\subsubsection{Centralizers over generic $\zeta$.}
When $\zeta$ is generic, the connected centralizer $J = J_{(x, \zeta)}$ is given by a slanted copy of the resonant subgroup $L = L_{[x, \zeta]}$.
\begin{lemma}\label{lemma: centralizer as resonant subgroup -- generic zeta}
Let $(x, \zeta) \in T \times \Gm^{\rot}$ with $\zeta$ generic. We have a preferred identification
\begin{equation*}
    \slant_{\omega}: J \xrightarrow{=} L,
\end{equation*}
describing the connected centralizer $J$ as a slanted copy of the resonant subgroup $L$. On root systems, this is induced by
$\slant_{\omega}: \Phi^{\bullet}_J \to \Phi^{\bullet}_L $, $\beta \mapsto \beta + m_{\beta}\hbar$.
\end{lemma}
\begin{proof}
We have described the root system of $J$ in Lemma \ref{lemma: root spaces of J}. Since $\ell = \ord(\zeta)=0$, it is given by $ \Phi^{\bullet}_J = \{ \beta + m_{\beta} \cdot \hbar \mid \beta \in \Phi^{\bullet}_L \}$ for the (uniquely determined) additive function $\mm: \beta \mapsto m_{\beta}$ from Notation \ref{notation: additive function m}.
Slanting by $\omega$ as in Lemma \ref{discussion: slanting J in general}, we get
$\Phi^{\bullet}_{\slant_{\omega}(J)} = \{ \beta \mid \beta \in \Phi^{\bullet}_L \}$, which is the root system of the connected reductive group $L$. The discussion in the same lemma implies the formula on root systems.
\end{proof}



\subsubsection{Labeling set for components.}
Under slanting on root systems $\slant_{\omega}: \Phi^{\bullet}_J \xrightarrow{=} \Phi^{\bullet}_L$ from Remark \ref{remark: slanting on weyl groups general case}, the group $W_J$ identifies as follows.
\begin{lemma}\label{lemma: resonant subgroup -- zeta generic}
Slanting identifies
\begin{equation*}
    \slant_{\omega}: W_J \xrightarrow{=} W_L.
\end{equation*}
\end{lemma}
\begin{proof}
Clear from Lemma \ref{lemma: centralizer as resonant subgroup -- generic zeta} since $W_J = W_{\Phi^{\bullet}_J}$ and $W_L = W_{\Phi^{\bullet}_L}$.   
\end{proof}


We can now describe the labeling set $\Theta$ from Notation \ref{notation: Theta} more explicitly.
To this end, recall the $\omega$-shifted action from \S \ref{section: shifted actions}.
\begin{lemma}\label{lemma: labeling set for components -- zeta generic}
For $(x, \zeta)$ with $\zeta$ generic, we have
\begin{equation*}
 \Theta = |\overline{\Gamma}_{(x, \zeta)}|  = X_{\bullet} / (W_{L}, \bullet_{\omega}).
\end{equation*}
\end{lemma}
\begin{proof}
 By Notation \ref{notation: Theta} and Lemma \ref{lemma: fiber of cocharacter graphs as lattice quotient}, we have $\Theta = |\overline{\Gamma}_{(x, \zeta)}|  = X_{\bullet}/W_J$ with respect to the natural action of $W_J$ on $X_{\bullet}$. Since the natural action of $W_J$ identifies with the $\omega$-shifted action of $\slant_{\omega}(W_J)$ by Remark \ref{remark: slanting on weyl groups general case}, we identify the quotient as
 \begin{equation*}
     X_{\bullet}/W_J = X_{\bullet}/(W_{\slant_{\omega}(J)}, \bullet_{\omega}).
 \end{equation*}
Since $W_{\slant_{\omega}(J)} = W_L$ is the resonant Weyl group by Lemma \ref{lemma: resonant subgroup -- zeta generic}, the result follows.
\end{proof}

\begin{remark}
Similar description works for other variants of cocharacter graphs -- for example,  
\begin{equation*}
   |{\Gamma}_{(x, \zeta)}|  = W_{\ext} / (W_{L}, \bullet_{\omega}). 
\end{equation*}
\end{remark}




\subsubsection{Positive stabilizers as parabolic subgroups.}
In terms of the slanted identification, the subgroup $J^+$ is given as $P_x$, the parabolic subgroup from \S \ref{section: parabolic and levi subgrops of L}. More generally, stabilizers of each $t^{\lambda} \in \Gr_G$ are described as follows.
\begin{lemma}\label{lemma: positive stabilizers as parabolics -- generic zeta}
For $(x, \zeta)$ with $\zeta$ generic and arbitrary $\lambda \in X^{\bullet}$, slanting identifies
\begin{equation*}
    \slant_{\omega}: J^{+}_{\lambda} \xrightarrow{=}  P_{x \cdot \lambda(\zeta)},
\end{equation*}   
where $P_{x \cdot \lambda(\zeta)} \leq L$ is a parabolic subgroup whose Levi factor $L_{x \cdot \lambda(\zeta)} \leq L$ is the connected stabilizer of $x \cdot \lambda(\zeta)$ inside $G$. 
\end{lemma}
\begin{proof}
Note that $U_{\beta + m\hbar} \subseteq t^{\lambda} \cdot \lop G \cdot t^{-\lambda}$ if and only if $ t^{-\lambda} \cdot U_{\beta + m\hbar} \cdot t^{\lambda} \subseteq \lop G$. Since $ t^{-\lambda} \cdot U_{\beta + m\hbar} \cdot t^{\lambda} = U_{\beta + (m - \langle \beta, \lambda \rangle)\hbar}$ by Lemma \ref{lemma: conjugation of affine root spaces}, the root spaces of $t^{\lambda} \cdot \lop G \cdot t^{-\lambda}$ are given by
$\{ \beta + m\hbar \mid m - \langle \beta, \lambda \rangle \geq 0 \}$.
%
To describe the roots spaces of $J^{+}_{\lambda}$, we intersect with the root system of $J$. This yields
\begin{equation*}
    \Phi^{\bullet}_{J^{+}_{\lambda}} = \{ \beta - m_{\beta}\hbar \mid \beta \in \Phi^{\bullet}_L, \  m_{\beta} + \langle \beta, \lambda \rangle \leq 0 \}
\end{equation*}
After slanting as in \S \ref{section: slanting}, this is the subgroup of $\slant_{\omega}(J^{+}_{\lambda}) \leq L$ spanned by the root spaces from
\begin{equation*}
    \Phi^{\bullet}_{\slant_{\omega}(J^{+}_{\lambda})} = \{ \beta \in \Phi^{\bullet}_L \mid  m_{\beta} + \langle \beta, \lambda \rangle \leq 0 \} = \{ \beta \mid \beta(x \cdot \lambda(\zeta) = \zeta^m \ \text{for some} \ m \leq 0 \}.
\end{equation*}
Here, the second equality follows since $\beta(x \cdot \lambda(\zeta)) = \beta(x) \cdot \beta(\lambda(\zeta)) = \zeta^{m_{\beta}} \cdot \zeta^{\langle \beta, \lambda \rangle} = \zeta^{m_{\beta} + \langle \beta, \lambda \rangle}$; this uniquely determines the exponent by genericity of $\zeta$.

These are the root spaces of $P_{x \cdot \lambda(\zeta)}$ from \eqref{equation: roots of Px}, a parabolic subgroup with Levi factor $L_{x \cdot \lambda(\zeta)}$ by Lemma \ref{lemma: parabolic and levi subgrops of L} and Remark \ref{remark: parabolic and levi subgrops of L}.
\end{proof}

\begin{remark}\label{remark: diferent labelings of centralizers -- generic zeta}
We may label the above subgroups of $L$ differently. Let $\vartheta \in \Theta = X_{\bullet}/(W_{\slant_{\omega}(J)}, \bullet_{\omega})$ be the orbit of $\lambda$. Under the interpretation, $\Theta = |\overline{\Gamma}_{(x, \zeta)}| \subseteq |T \sslash W|$, the element $\vartheta$ is given by the equivalence class of the point $x \cdot \lambda(\zeta)$. In Notation \ref{notation: levi and paraolic subgroups labeled by elements of Theta}, we get
\begin{equation*}
  P_{\vartheta} = P_{x \cdot \lambda(\zeta)} = P_{\omega+\lambda^{\ad}} 
  \qquad \text{and} \qquad
  L_{\vartheta} = L_{x \cdot \lambda(\zeta)} = L_{\omega+\lambda^{\ad}},
\end{equation*}
where the latter equalities realize our subgroups as the parabolic and Levi subgroup of $L$ attached to the cocharacter $\omega + \lambda^{\ad} \in X^{\ad}_{\bullet, L}$ by Remark \ref{remark: parabolic and levi subgrops of L}.
\end{remark}


\begin{lemma}\label{lemma: positive centralizer as parabolic -- gneric zeta}
Inside $\lop G \rtimes \Gm^{\rot}$,  
\begin{equation*}
\slant_{\omega}: J^{+} \xrightarrow{=} P_x 
\end{equation*}
where $P_x = P_{\omega} \leq L$ is the opposite parabolic subgroup with Levi factor $L_x = L_{\omega} \leq L$.
\end{lemma}
\begin{proof}
Special case of Lemma \ref{lemma: positive stabilizers as parabolics -- generic zeta} and Remark \ref{remark: diferent labelings of centralizers -- generic zeta} for $\lambda$ trivial.    
\end{proof}



\subsubsection{Fixed points on affine Grassmannian at \texorpdfstring{$\zeta$}{zeta} generic.}
We now describe the fixed points on the whole affine Grassmannian.
\begin{theorem}\label{lemma: fixed points of (c, zeta) on Gr_G -- zeta generic}
For $(x, \zeta) \in T \times \Gm^{\rot}$ with $\zeta$ generic, we have
\begin{equation}\label{equation: fixed points of (x, zeta) on Gr_G -- zeta generic}
    \Fix^{\red}_{(x, \zeta)}(\Gr_G) = \coprod_{\vartheta \in X_{\bullet}/(W_{L}, \bullet_{\omega})} L/P_{\vartheta}
\end{equation}
where $L = L_{[x, \zeta]}$ is the resonant subgroup, $P_{\vartheta}$ is the parabolic subgroup labeled by $\vartheta \equiv x \cdot \lambda(\zeta)$, and $W_L$ acts on $X_{\bullet}$ by the $\omega$-shifted action.
\end{theorem}
\begin{proof}
By Lemma \ref{lemma: fixed points on affine Grassmannian}, the fixed points are given by a disjoint union of the homogeneous spaces $J/J^+_{\vartheta}$ ranging over $\vartheta \in \Theta$.
In the situation at hand, $\Phi^{\bullet}_J \cong \Phi^{\bullet}_L$ and $J \cong L$ by the slanted identification from Lemma \ref{lemma: centralizer as resonant subgroup -- generic zeta}. The labeling set for components from Notation \ref{notation: Theta} becomes $\Theta \cong X_{\bullet} / (W_L, \bullet_{\omega})$ by Lemma \ref{lemma: labeling set for components -- zeta generic}. For given $\vartheta \in X_{\bullet} / (W_L, \bullet_{\omega})$, pick a lift $\lambda \in X_{\bullet}$. Slanting identifies $J^+_{\lambda} \cong P_{x \cdot \lambda(\zeta)}$ by Lemma \ref{lemma: positive stabilizers as parabolics -- generic zeta}; this parabolic is independent of the lift and denoted $P_{\vartheta}$ as in Notation \ref{notation: levi and paraolic subgroups labeled by elements of Theta} and Remark \ref{remark: diferent labelings of centralizers -- generic zeta}. The result follows.
\end{proof}




\begin{remark}
The above theorem can be also proved through the viewpoint of \S \ref{section: fixed points on GrG and GrL}, which we outline now. Passing to appropriate components if necessary, we only need to compute fixed points of the loop rotation by $\zeta$ on $\Gr_{L^{\ad}}$. These fixed-points are given by
\begin{equation}
    \Fix^{\red}_{(1, \zeta)}(\Gr_{L^{\ad}}) \cong \coprod_{\lambda \in X^{\ad}_{\bullet, L}/W_L} L^{\ad}/P_{\lambda}.
\end{equation}
where $\lambda$ runs over orbit representatives for the usual action of $W_L$ on $X^{\ad}_{\bullet, L}$ and $P_{\lambda} \leq L^{\ad}$ is the parabolic subgroup associated to $\lambda$ as in \S \ref{section: levi subgroups}.

Indeed, since $\omega$ is a cocharacter of $L^{\ad}$, shifting identifies $X^{\ad}_{\bullet, L}/(W_L, \bullet_{\omega}) \cong X^{\ad}_{\bullet, L}/W_L$. The orbits of this naive action $W_L \acts X^{\ad}_{\bullet, L}$ parametrize the dominant cocharacters $X^+_{\bullet}(L^{\ad})$ of $L^{\ad}$. 
The fixed points of $(1, \zeta)$ on any affine Grassmannian are well-known, e.g. \cite[above Example 2.1.12]{Zhu15}. Each affine Schubert cell is a geometric vector bundle over the corresponding partial flag variety, with nontrivial weights in the fibers. The loop rotation fixed points are then given by the zero section. These are the same as the fixed points of the single element $\zeta$ by the genericity assumption.
\end{remark}



\subsubsection{Fixed points on affine Schubert varieties as finite Schubert varieties.}
We can now prove the desired description of fixed-points on any affine Schubert variety.
\begin{theorem}\label{theorem: generic zeta -- fiber of fixed points and classical schubert varieties}
The reduced scheme $\Fix^{\red}_{(x, \zeta)}(\Gr_{\leq \mu})$ is isomorphic to a union of classical Schubert varieties for the connected reductive group $L = L_{[x, \zeta]}$.
\end{theorem}
\begin{proof}
Unraveling Lemma \ref{lemma: upshot of the general case} via the slanted descriptions from Lemmata \ref{lemma: positive centralizer as parabolic -- gneric zeta},  \ref{lemma: fixed points of (c, zeta) on Gr_G -- zeta generic} identifies the variety 
$F = \Fix^{\red}_{(x, \zeta)}(\Gr_{\leq \mu})$
of our interest as a closed $P_{x}$-equivariant subvariety
\begin{equation*}\label{equation: B-equivariant closed immersion into classical flag varieties -- zeta generic}
P_{x} \acts F \hookrightarrow \coprod_{\vartheta \in X_{\bullet}/(W_{L}, \bullet_{\omega})} L/P_{\vartheta}   
\end{equation*}
where $P_{x} \leq L$ acts by left translation on each of the partial flag varieties $L/P_{\vartheta} = L/P_{x \cdot \lambda(\zeta)}$

We may work on each connected components $L/P_{\vartheta}$ separately. The corresponding piece of the fixed point scheme $F^{\vartheta} = F \cap L/P_{\vartheta}$ cuts out a closed $P_{x}$-equivariant subvariety -- therefore, it is a downwards closed union of finitely many Schubert cells. Altogether, we get a (not necessarily disjoint) union of classical Schubert varieties in $L/P_{\vartheta}$.
\end{proof}


\subsubsection{Combinatorial description.}
Moreover, the classical Schubert varieties appearing in Theorem \ref{theorem: generic zeta -- fiber of fixed points and classical schubert varieties} can be described combinatorially. Indeed, to get such a description, all we need to do is to describe the set of relevant torus fixed-points. 
%
These fixed points on $\Gr_{\leq \mu}$ are labeled by the finite set $\wt(\mu)$ from \eqref{notation: weights in mu}.

\begin{lemma}\label{lemma: combinatorial description -- zeta generic}
The $T\times \Gm^{\rot}$-fixed-points on $\Fix^{\red}_{(x, \zeta)}(\Gr_{\leq \mu})$ are given by $\{ t^{\lambda} \mid \lambda \in \wt(\mu) \}$. Two such fixed-points $t^{\lambda_1}$, $t^{\lambda_2}$ lie in the same connected component of $\Fix^{\red}_{(x, \zeta)}(\Gr_{\leq \mu})$ if and only if $\lambda_1$, $\lambda_2$ lie in the same orbit of the $\omega$-shifted action $(W_L, \bullet_{\omega}) \acts X_{\bullet}$.    
\end{lemma}
\begin{proof}
Follows by unraveling Lemma \ref{lemma: connected components of fixed points n affine schubert variety}.
Under these identifications, wherever $\lambda_1$ and $\lambda_2$ are related by a reflection from $(W_L, \bullet_{\omega})$, the two points $t^{\lambda_1}$, $t^{\lambda_2}$ can be connected by a $\P^1$ inside $F$.
\end{proof}

\begin{remark}
For an alternative proof of Lemma \ref{lemma: combinatorial description -- zeta generic}, we only need to see that for each $\vartheta \in \Theta$, the subscheme $F^{\vartheta}= \Fix^{\red}_{(x, \zeta)}(\Gr_{\leq \mu}) \cap L/P_{\vartheta}$ is already connected. This is clear since $F^{\vartheta}$ is a closed union of Bruhat cells by Lemma \ref{theorem: generic zeta -- fiber of fixed points and classical schubert varieties} and the finite Bruhat order has a unique minimal element.
\end{remark}


\begin{remark}[Reducible unions]
By the above result, the connected components of the fixed points are given by $\Fix^{\red, \vartheta}_{(x, \zeta)}(\Gr_{\leq \mu}) \subseteq  L/P_{\vartheta}$, $\vartheta \in \Theta$; each of those is a union of finite Schubert varieties. However, it may not be irreducible. 

An explicit example already appears when $G= GL_4$, $\Gr_{\leq \mu}$ for $\mu = (4,4,0,0)$, generic $\zeta$, and $x = (\zeta^3, \zeta^2, \zeta^1, 1)$. Then $\omega = \rho$, so the $\omega$-shifted action corresponds to the dot action. In total there are $|\wt(\mu)| = 85$ fixed points. The partition into connected components of $\Fix^{\red}_{(x, \zeta)}(\Gr_{\leq \mu})$ is given by
\begin{multline*}
    85 
    = 14 + 8 + 4 + 2 + 1 + 4 
    + 6 + 3 + (4+3-2) + (3+2-1)\\
    + 3 + (3+2-1) +3 + (2+4-1)
    +(4+3-2) + 2 + 1 + 1
    + 3 + 2 + 2 + 1
    +2
\end{multline*}
Here, each bracketed term corresponds to a single connected component given by a reducible union of two finite Schubert varieties intersecting along a third one.
\end{remark}



\subsection{\texorpdfstring{$\zeta$}{Zeta} root of unity}\label{section: zeta root of unity}

\subsubsection{Setup.}
We now discuss the complementary case when $\zeta$ is a root of unity. Fix such point $(x, \zeta) \in T \times \Gm^{\rot}$ and denote by $\ell = \ord(\zeta) \in \N$ the primitive order of $\zeta$.
%
We describe the fiber $\Fix^{\red}_{(x, \zeta)}(\Gr_G)$ as a finite union of $\ell$-dilated affine flag varieties for the resonant subgroup $L = L_{[x, \zeta]}$. Inside them, $\Fix^{\red}_{(x, \zeta)}(\Gr_{\leq \mu})$ is given by a closed union of $\ell$-dilated Iwahori orbits.



\subsubsection{Centralizers over roots of unity.}
\begin{lemma}\label{lemma: zeta root of unity -- computation of centralizer}
Inside $\lo L$, slanting identifies
\begin{equation*}
    \slant_{\omega}: J \xrightarrow{=} \lol L^{\sc}.
\end{equation*}
\end{lemma}
\begin{proof}
By Lemma \ref{lemma: root spaces of J}, the root subsystem of $J$ is given by  
$\Phi^{\bullet}_J = \{\beta - (m_{\beta} + \ell \Z)\hbar \mid \beta \in \Phi^{\bullet}_{L}\}$.
Slanting by $\omega$ as in Discussion \ref{discussion: slanting J in general}, we obtain
\begin{equation*}
\Phi^{\bullet}_{\slant_{\omega}(J)} = \{\beta - \ell \Z\hbar \mid \beta \in \Phi^{\bullet}_{L}\},
\end{equation*}
which is precisely the root subsystem $\Phi^{\bullet}_{\lol L^{\sc}}$ of $\lol L^{\sc} = \lol^{\circ} L$, embedded in $\lo G$ naively. Since both $J$ and $\lol L^{\sc}$ are connected, the result follows.
\end{proof}

\begin{remark}
As a sanity check, note that an element $t^{\lambda}$ in $\lo T$ lies in the centralizer of $(x, \zeta)$ if and only if $\lambda(\zeta) = 1$, in other words if and only if $\lambda \in \ell X_{\bullet}$. Consequently, $\Cent_{\lo G}(x) \cap \lo T = \lol T$. The identity componet of the centralizer then intersects $\lo T$ as $J \cap \lo T = \lol T^{\sc} = \slant_{\omega}(\lol T^{\sc})$.
\end{remark}



\subsubsection{Labeling set for components.}
Slanting on root systems $\slant_{\omega}: \Phi^{\bullet}_J \xrightarrow{=} \Phi^{\bullet}_{\lol L^{\sc}}$ from Remark \ref{remark: slanting on weyl groups general case} identifies the corresponding reflection subgroups.
\begin{lemma}\label{lemma: resonant subgroup -- zeta root of unity}
Slanting identifies
\begin{equation*}
    \slant_{\omega}: W_J \xrightarrow{=} \ell \Z\Phi^{\bullet}_L \rtimes W_L.
\end{equation*}
as the $\ell$-dilated affine Weyl group of $L$.
\end{lemma}
\begin{proof}
Clear from Lemma \ref{lemma: zeta root of unity -- computation of centralizer} since $W_J = W_{\Phi^{\bullet}_J}$ and $\ell \Z\Phi^{\bullet}_L \rtimes W_L = W_{\Phi^{\bullet}_{\lol L^{\sc}}}$.   
\end{proof}


We can now describe the labeling set $\Theta$ from Notation \ref{notation: Theta} as follows.
\begin{lemma}\label{lemma: labeling set for components -- zeta root of unity}
For $(x, \zeta)$ with $\zeta$ generic, we have
\begin{equation*}
 \Theta = |\overline{\Gamma}_{(x, \zeta)}| = X_{\bullet} / (W^{\aff}_{L}, \bullet^{\ell}_{\omega}).
\end{equation*}
\end{lemma}
\begin{proof}
 By Notation \ref{notation: Theta} and Lemma \ref{lemma: fiber of cocharacter graphs as lattice quotient}, we have $\Theta = |\overline{\Gamma}_{(x, \zeta)}|  = X_{\bullet}/W_J$ with respect to the natural action of $W_J$ on $X_{\bullet}$. 
%
Since the natural action of $W_J$ identifies with the $\omega$-shifted action of $\slant_{\omega}(W_J)$ by Remark \ref{remark: slanting on weyl groups general case}, we obtain the first identification in
 \begin{equation*}
     X_{\bullet}/W_J 
     = X_{\bullet}/(W_{\slant_{\omega}(J)}, \bullet_{\omega}) 
     = X_{\bullet}/(W_{L, \aff}, \bullet^{\ell}_{\omega}). 
 \end{equation*}
 The second identification follows since $W_{\slant_{\omega}(J)} = \Z\ell\Phi^{\bullet}_L \rtimes W_L$ is the $\ell$-dilated affine Weyl group of $L$ by Lemma \ref{lemma: resonant subgroup -- zeta root of unity}, whose orbits on $X_{\bullet}$ identify with the orbits of $W_{L, \aff}$ under the $\ell$-dilated action.
\end{proof}

\begin{remark}
Similar description works for other variants of cocharacter graphs -- for example,  
\begin{equation*}
   |{\Gamma}_{(x, \zeta)}|  = W_{\ext} / (W_{L, \aff}, \bullet^{\ell}_{\omega}). 
\end{equation*}
\end{remark}



\subsubsection{Positive centralizers as parahoric subgroups.}
Under slanting, we identify the positive centralizer $J^+$ as the $\ell$-dilated parahoric subgroup $\eP^{x}_{\ell} = \eP^{\omega}_{\ell}$ of $\lol L^{\sc}$.
More generally, we describe the stabilizers $J^{+}_{\lambda}$, $\lambda \in X_{\bullet}$.

\begin{lemma}\label{lemma: zeta root of unity -- positive stabilizers}
Inside $\lop G$, slanting identifies
\begin{equation*}
\slant_{\omega}: J^{+}_{\lambda} \xrightarrow{=} \eP^{x \cdot \lambda(\zeta)}_{\ell},
\end{equation*}    
where $\eP^{x \cdot \lambda(\zeta)}_{\ell} = \eP_{\ell}^{\omega+\lambda^{\ad}}$ is the standard opposite parahoric subgroup of $\lol L^{\sc}$ labeled by the point $x \cdot \lambda(\zeta) \in T$ resp. the cocharacter $\omega+\lambda^{\ad} \in X^{\ad}_{\bullet, L}$.
\end{lemma}
\begin{proof}
Note that $U_{\beta + m\hbar} \subseteq t^{\lambda} \cdot \lop G \cdot t^{-\lambda}$ if and only if $ t^{-\lambda} \cdot U_{\beta + m\hbar} \cdot t^{\lambda} \subseteq \lop G$. Since $ t^{-\lambda} \cdot U_{\beta + m\hbar} \cdot t^{\lambda} = U_{\beta + (m - \langle \beta, \lambda \rangle)\hbar}$ by Lemma \ref{lemma: conjugation of affine root spaces}, the root spaces of $t^{\lambda} \cdot \lop G \cdot t^{-\lambda}$ are given by
$\{ \beta + m\hbar \mid m - \langle \beta, \lambda \rangle \geq 0 \}$.
%
To describe the roots of $J^{+}_{\lambda}$, we intersect with the root system of $J$. This yields
\begin{equation*}
    \Phi^{\bullet}_{J^{+}_{\lambda}} = \{ \beta - (m_{\beta} + c\ell)\hbar \mid \beta \in \Phi^{\bullet}_L, \ c \in \Z, \ m_{\beta} + c\ell + \langle \beta, \lambda \rangle \leq 0 \}
\end{equation*}
Slanting by $\omega$ as in \S \ref{section: slanting} and rewriting $m_{\beta} + \langle \beta, \lambda \rangle = \langle \beta, \omega \rangle + \langle \beta, \lambda^{\ad} \rangle = \langle \beta, \omega + \lambda^{\ad} \rangle$, this becomes
\begin{equation*}
    \Phi^{\bullet}_{\slant_{\omega}(J^{+}_{\lambda})} = \{ \beta - c\ell\hbar \mid \beta \in \Phi^{\bullet}_L, \ c \in \Z, \ c\ell + \langle \beta, \omega+\lambda \rangle \leq 0 \}
\end{equation*}
where $\lambda^{\ad} \in X^{\ad}_{\bullet, L}$ is the projection of $\lambda \in X_{\bullet}$.

These are precisely the root spaces of the standard opposite parahoric subgroup $\eP^{\omega+\lambda^{\ad}}_{\ell}$ of the $\ell$-dilated loop group $\lol L^{\sc}$ associated to the cocharacter $\omega+\lambda^{\ad} \in X_{\bullet, L}^{\ad}$ as in \S \ref{subsubsection: parahoric subgroups}, \S \ref{subsubsection: dilated variants}. 
\end{proof}

\begin{remark}\label{remark: positive stabilizers -- zeta root of unity}
By Remark \ref{remark: independece of J vartheta on the lift}, the stabilizer $J^{+}_{\lambda}$ depends, up to isomorphism given by conjugation, only on the class $\vartheta \in \Theta = X_{\bullet} / (W^{\aff}_{L}, \bullet^{\ell}_{\omega})$. Denote $\vartheta$ the image of $\lambda$ under the projection to $\Theta$. In terms of the identification $\Theta = |\overline{\Gamma}_{(x, \zeta)}| \subseteq |T \sslash W|$, the element $\vartheta$ is given by the equivalence class of the point $x \cdot \lambda(\zeta)$. We denote the parahoric subgroup from Lemma \ref{lemma: zeta root of unity -- positive stabilizers} by
\begin{equation*}
 \eP^{\vartheta}_{\ell} 
 = \eP^{x \cdot \lambda(\zeta)}_{\ell} 
 = \eP_{\ell}^{\omega+\lambda^{\ad}}.  
\end{equation*}    
\end{remark}

In particular, we deduce the description of the positive centralizer.
\begin{lemma}\label{lemma: zeta root of unity -- positive centralizer}
\label{lemma: H^+ conjugates to P via alpha general case roots of unity}
Inside $\lop L$, slanting identifies
\begin{equation*}
\slant_{\omega}: J^{+} \xrightarrow{=} \eP_{\ell}^{x}.
\end{equation*}
where $\eP_{\ell}^{x} = \eP_{\ell, L}^{\omega} \leq \lo_{\ell} L$ is the standard opposite parahoric subgroup corresponding to $\omega$.
\end{lemma}
\begin{proof} 
Special case of Lemma \ref{lemma: zeta root of unity -- positive stabilizers} for trivial $\lambda$.
\end{proof}




\subsubsection{Fixed points on affine Grassmannian at \texorpdfstring{$\zeta$}{zeta} root of unity.}

We can now make explicit the fixed points on the whole affine Grassmannian.
\begin{theorem}\label{lemma: fixed points of (c, zeta) on Gr_G -- zeta root of unity}
Let $(x, \zeta) \in T \times \Gm^{\rot}$, where $\zeta$ is a root of unity of $\ord(\zeta) = \ell$. Then
\begin{equation*}
    \Fix^{\red}_{(x, \zeta)}(\Gr_G) = \coprod_{\vartheta \in X_{\bullet}/(W_{L, \aff}, \bullet^{\ell}_{\omega})} \Fl^{\vartheta}_{\ell, L^{\sc}}.
\end{equation*}
where $L=L_{[x, \zeta]}$ is the resonant subgroup, $\Fl^{\vartheta}_{\ell, L^{\sc}} = \lol L^{\sc}/ \eP^{\vartheta}_{\ell}$ is the $\ell$-dilated affine flag variety labeled by $\vartheta \equiv x \cdot \lambda(\zeta)$, and $W_{L, \aff}$ acts on $X_{\bullet}$ by the $\omega$-shifted $\ell$-dilated action.
\end{theorem}
\begin{proof}
Unravel Lemma \ref{lemma: fixed points on affine Grassmannian} in terms of the explicit identification above. The fixed points are a disjoint union of homogeneous spaces $J/J^+_{\vartheta}$, $\vartheta \in \Theta$. We identify $J = \lol L^{\sc}$ by Lemma \ref{lemma: zeta root of unity -- computation of centralizer}. The labeling set for components from Notation \ref{notation: Theta} becomes $\Theta = X_{\bullet} / \Theta = X_{\bullet}/(W_{\aff}, \bullet^{\ell}_{\omega} ) $ by Lemma \ref{lemma: labeling set for components -- zeta root of unity}. 
For given $\vartheta$, pick a lift $\lambda \in X_{\bullet}$. Slanting identifies $J^+_{\vartheta} = \eP^{\vartheta}_{\ell} = \eP^{x \cdot \lambda(\zeta)}_{\ell}$ by Lemma \ref{lemma: zeta root of unity -- positive stabilizers} and Remark \ref{remark: positive stabilizers -- zeta root of unity}. The result follows.
\end{proof}


\begin{remark}
For an alternative approach to the above theorem, we may again adapt the viewpoint of \S \ref{section: fixed points on GrG and GrL}, reducing the problem to the computation of (appropriate components) of the loop rotation by $\zeta$ on $\Gr_{L^{\ad}}$.

Since $\zeta$ is a primitive root of unity, these are given by
\begin{equation}\label{equation: alternative approach roots of unity}
    \Fix^{\red}_{(1, \zeta)}(\Gr_{L^{\ad}}) = \coprod_{\vartheta \in X^{\ad}_{\bullet, L}/ (W_{L, \aff}, \bullet^{\ell}_0)} \lol L^{\sc}/\eP^{\lambda}_{\ell}
\end{equation}
where $\lambda$ runs over the orbits of $(W_{\aff}, \bullet^{\ell}_0) \acts X_{\bullet}$ and $\eP^{\lambda}_{\ell} \leq \lol L^{\sc}$ is the associated standard opposite parahoric subgroup of the $\ell$-dilated loop group.

Here, we have identified $ X_{\bullet}/ (W_{\aff}, \bullet^{\ell}_{\omega}) \cong X_{\bullet}/ (W_{\aff}, \bullet^{\ell}_{0})$ by shifting with $\omega$, which is a cocharacter of $L^{\ad}$. 
Since $\ord(\zeta) = \ell$, the group $\mu_{\ell}$ is a finite cyclic group of order $\ell$ on the generator $\zeta$, so \eqref{equation: alternative approach roots of unity} follows by \cite[\S 1.6, \S 4.6]{RW22} or \cite[\S 2.1.4]{BBSV22}.
\end{remark}


\subsubsection{Fixed points on affine Schubert varieties as parahoric orbits.}
We now intersect the above description with any affine Schubert variety $\Gr_{\leq \mu}$ of $\Gr_G$.
\begin{theorem}\label{theorem: fiber of fixed points and classical schubert varieties general case at root of unity}
The scheme $\Fix^{\red}_{(x, \zeta)}(\Gr_{\leq \mu})$ is isomorphic to a closed union of $\eP_{\ell, L^{\sc}}^{\omega}$-orbits -- in particular $\Iw_{\ell, L^{\sc}}$-orbits -- on affine flag varieties $\Fl_{\ell, L}^{\vartheta}$ for the dilated loop group $\lo_{\ell} L^{\sc}$ of the simply connected form of the resonant subgroup $L = L_{[x, \zeta]}$.
\end{theorem}
\begin{proof}
Unraveling Lemma \ref{lemma: upshot of the general case} via the slanted descriptions from Lemmata \ref{lemma: fixed points of (c, zeta) on Gr_G -- zeta root of unity} and \ref{lemma: zeta root of unity -- positive centralizer} identifies the reduced scheme
$F = \Fix^{\red}_{(x, \zeta)}(\Gr_{\leq \mu})$
of our interest as a $\eP_{\ell}^{\omega}$-equivariant subvariety
\begin{equation*}\label{equation: B-equivariant closed immersion into classical flag varieties -- zeta rot of unity}
\eP_{\ell}^{\omega} \acts F \hookrightarrow \coprod_{\vartheta \in X_{\bullet}/(W_{L, \aff}, \bullet^{\ell}_{\omega})} \lo_{\ell} L^{\sc}/\eP_{\ell}^{\vartheta}   
\end{equation*}
where $\eP_{\ell}^{\omega} \leq \lo_{\ell} L^{\sc}$ acts by left translation on each of the affine flag varieties $\Fl_{\ell, L^{\sc}}^{\vartheta} = \lol L^{\sc}/\eP_{\ell}^{\vartheta}$.

Restricting attention to a single component $\Fl_{\ell, L^{\sc}}^{\vartheta}$ of the right-hand side, we get a closed union of $\eP_{\ell}^{\omega}$-orbits as desired. Since $\Iw_{\ell, L^{\sc}} \leq \eP_{\ell}^{\omega}$, it is in particular closed union of $\ell$-dilated Iwahori orbits.
\end{proof}


\subsubsection{Combinatorial description.}
Again, the dilated Iwahori orbits appearing in Theorem \ref{theorem: fiber of fixed points and classical schubert varieties general case at root of unity} may be described via the relevant torus fixed-points $t^{\lambda}$, labeled by the finite set $\wt(\mu)$ from \eqref{notation: weights in mu}.

\begin{lemma}\label{lemma: combinatorial description -- zeta root of unity}
The $T\times \Gm^{\rot}$-fixed-points on $\Fix^{\red}_{(x, \zeta)}(\Gr_{\leq \mu})$ are given by $\{ t^{\lambda} \mid \lambda \in \wt(\mu) \}$. Two such fixed-points $t^{\lambda_1}$, $t^{\lambda_2}$ lie in the same connected component of $\Fix^{\red}_{(x, \zeta)}(\Gr_{\leq \mu})$ if and only if $\lambda_1$, $\lambda_2$ lie in the same orbit of the $\omega$-shifted $\ell$-dilated action $(W_{L, \aff}, \bullet^{\ell}_{\omega}) \acts X_{\bullet}$.    
\end{lemma}
\begin{proof}
Follows by unraveling Lemma \ref{lemma: connected components of fixed points n affine schubert variety} in terms of Lemma \ref{lemma: zeta root of unity -- positive centralizer} Theorem \ref{theorem: fiber of fixed points and classical schubert varieties general case at root of unity}.
\end{proof}

As in the finite case, Lemma \ref{lemma: combinatorial description -- zeta root of unity} shows that each $F^{\vartheta} = \Fix^{\red}_{(x, \zeta)}(\Gr_{\leq \mu}) \cap \lol L^{\sc}/\eP_{\ell}^{\vartheta}$ is already connected.